\section*{Sets and Indices}

\[
V=\{0,1,\dots,n\}
\]
set of all nodes, where node \(0\) is the depot and nodes \(1,\dots,n\) are customers.

\[
N=\{1,\dots,n\}
\]
set of customers.

\[
A \subseteq \{(i,j)\in V\times V : i\neq j\}
\]
set of directed arcs retained by the preprocessing in the code. An arc \((i,j)\) is included only if it passes the following feasibility screen:
\[
(i,j)\in A \iff i\neq j \text{ and}
\]
\[
\begin{cases}
e_i+s_i+d_{ij}\le l_j, & i\in N,\ j\in N,\\
e_0+d_{0j}\le l_j, & i=0,\ j\in N,\\
e_i+s_i+d_{i0}\le l_0, & i\in N,\ j=0.
\end{cases}
\]

Indices:
\[
i,j \in V.
\]

\section*{Parameters}

\[
n=20
\]
number of customers (\texttt{num\_customers}).

\[
K^{\max}=5
\]
maximum number of trucks (\texttt{max\_trucks}).

\[
Q=200
\]
truck capacity (\texttt{truck\_capacity}).

\[
(x_i,y_i)
\]
coordinates of node \(i\) (\texttt{coords[i]}), with depot at \((40,50)\).

\[
q_i
\]
demand of customer \(i\) (\texttt{demands\_list[i]}), with \(q_0=0\).

\[
[e_i,l_i]
\]
time window of node \(i\) (\texttt{tw\_list[i]}). In particular, the depot has
\[
[e_0,l_0]=[0,1236].
\]

\[
s_i
\]
service duration at node \(i\) (\texttt{svc\_list[i]}), with \(s_0=0\). In the data, customer service times are \(90\) minutes.

\[
d_{ij}=\sqrt{(x_i-x_j)^2+(y_i-y_j)^2}
\]
travel distance (and, in the implemented timing constraints, also travel time) from node \(i\) to node \(j\) (\texttt{dist[i,j]}).

\[
M=l_0+\max_{i\in V} s_i
\]
big-\(M\) constant used in the time propagation constraints (\texttt{M\_val}).

\section*{Decision Variables}

\[
x_{ij} \in \{0,1\} \qquad \forall (i,j)\in A
\]
binary arc-selection variable indicating whether arc \((i,j)\) is used (\texttt{x[i,j]}).

\[
t_i \ge 0 \qquad \forall i\in V
\]
service start time at node \(i\) (\texttt{t[i]}).

\section*{Objective}

Minimize total traveled distance:
\[
\min \sum_{(i,j)\in A} d_{ij}\,x_{ij}.
\]

\section*{Constraints}

Depot departure limit:
\[
\sum_{j\in N:\,(0,j)\in A} x_{0j} \le K^{\max}.
\]

Exactly one incoming arc for each customer:
\[
\sum_{j\in V:\,(j,i)\in A} x_{ji}=1
\qquad \forall i\in N.
\]

Exactly one outgoing arc for each customer:
\[
\sum_{j\in V:\,(i,j)\in A} x_{ij}=1
\qquad \forall i\in N.
\]

Customer time-window bounds:
\[
t_i \ge e_i \qquad \forall i\in N,
\]
\[
t_i \le l_i \qquad \forall i\in N.
\]

Depot time-window bounds:
\[
t_0 \ge e_0,
\]
\[
t_0 \le l_0.
\]

Time propagation on used arcs ending at a customer:
\[
t_i + s_i + d_{ij} - t_j \le M(1-x_{ij})
\qquad \forall (i,j)\in A \text{ with } j\neq 0.
\]

Pairwise capacity screening on customer-to-customer arcs:
\[
q_i+q_j \le Q + Q(1-x_{ij})
\qquad \forall (i,j)\in A \text{ with } i\neq 0,\ j\neq 0.
\]

Binary and nonnegativity restrictions:
\[
x_{ij}\in\{0,1\} \qquad \forall (i,j)\in A,
\]
\[
t_i\ge 0 \qquad \forall i\in V.
\]

\section*{Notes on Implementation}

This formulation matches the implemented model exactly.

The code builds the arc set \(A\) by removing arcs that are impossible with respect to the destination time-window upper bound using earliest possible departure from the tail node. This preprocessing is part of the implemented model.

No explicit vehicle index is used. Routes are represented implicitly by arcs leaving and returning to the depot. The only fleet-size restriction is the bound on the number of arcs departing the depot.

There is no full cumulative vehicle-capacity formulation in the code. Instead, the only capacity-related constraint is the pairwise condition on any selected customer-to-customer arc:
\[
x_{ij}=1 \implies q_i+q_j \le Q.
\]
Thus the implemented model does not track total load along an entire route.

There are also no explicit subtour-elimination constraints. Time propagation constraints prevent cycles among customers that would be temporally inconsistent, but the formulation is not strengthened beyond the exact coded constraints.

The variable \(t_0\) is given an initial value \(0\) in the code (\texttt{t[0].setInitialValue(0)}), but this is only a solver hint and not an additional constraint.

The model is solved as a mixed-integer linear program using Gurobi through the PuLP-compatible interface (\texttt{GUROBI\_CMD}) with a \(100\)-second time limit.
